% Source: https://artofproblemsolving.com/community/c3h3393379_find_the_domain_and_range_of_the_following_function_sqrt4xfrac1sqrt
% Stable identity: aops-topic-3393379-q01
The denominator requires $9-x^2>0$, hence the domain is $\boxed{(-3,3)}$; this also ensures $4-x>0$. Differentiating gives $f'(x)=-\dfrac{1}{2\sqrt{4-x}}+\dfrac{x}{(9-x^2)^{\dfrac{3}{2}}}$. This is negative for $x\le0$. For $0<x<3$, its sign equals that of $H(x)-1$, where $H(x)=\dfrac{2x\sqrt{4-x}}{(9-x^2)^{\dfrac{3}{2}}}$. Its logarithmic derivative has positive denominator $2x(4-x)(9-x^2)$ and numerator $72-27x+16x^2-3x^3$. Since $0<x<3$, the numerator is greater than $72-27x+7x^2>0$; the last quadratic has negative discriminant and positive leading coefficient. Thus $H$ strictly increases from $0$ to $+\infty$. There is a unique $\alpha\in(0,3)$ with $H(\alpha)=1$, equivalently $4\alpha^2(4-\alpha)=(9-\alpha^2)^3$. The function decreases to $f(\alpha)$ and increases afterward, tending to infinity at both domain endpoints. Its range is therefore $\boxed{[\sqrt{4-\alpha}+\dfrac{1}{\sqrt{9-\alpha^2}},\infty)}$, where $\alpha\approx2.382445466$ and $f(\alpha)\approx1.820325985$.
